\section{Setting: cells, rooftops and moments}
\label{sec:setting}

\defTitle{Uniform grid}{grid} Let $a,b>0$ be the cell sizes in the $x$ and $y$ directions. For
$(m,n)\in\Zn^2$ the cell is
\begin{equation*}
    \Om{mn} = [ma,(m+1)a]\times[nb,(n+1)b] \subset \Rn^2 .
\end{equation*}
A point $\V{r}=(x,y)\in\Om{mn}$ has the local coordinates
\begin{equation*}
    \xi(\V{r}) = \frac{x}{a}-m \in[0,1], \qquad \eta(\V{r}) = \frac{y}{b}-n\in[0,1].
\end{equation*}
For an offset $\offs=(d_x,d_y)\in\Zn^2$ we write $\Om{\offs}=\Om{d_xd_y}$. The two cells $\Om{\V{0}}$ and $\Om{\offs}$ share
at least one point if and only if $\offs\in\mathcal{T}=\{-1,0,1\}^2$. We call these offsets touching
(Figure~\ref{fig:touching}): $\offs=\V{0}$ is the coincident cell, the four offsets $(\pm1,0)$, $(0,\pm1)$ share an edge
and the four offsets $(\pm1,\pm1)$ a vertex.

\defTitle{Rooftop functions}{rooftop} The $x$-directed rooftop $\V{f}^x_{mn}$ is supported on
$\Om{mn}\cup\Om{m+1,n}$ and
\begin{equation*}
    \V{f}^x_{mn}(\V{r}) = \UV{x}\cdot
    \begin{cases}
        \xi(\V{r}), & \V{r}\in\Om{mn} \quad\text{(rising wedge)},\\
        1-\xi(\V{r}), & \V{r}\in\Om{m+1,n} \quad\text{(falling wedge)}.
    \end{cases}
\end{equation*}
Its normal component on the common edge is one, and its divergence is $+1/a$ on $\Om{mn}$ and $-1/a$ on $\Om{m+1,n}$.
The $y$-directed rooftop $\V{f}^y_{mn}$ is defined in the same way with $\eta$, $b$ and $\Om{m,n+1}$. These are the rooftops
of~\cite{machacek2025method} and of \src{} (unit height on the edge).

\begin{figure}[ht]
    \centering
    % The touching neighbourhood of the reference cell and the x-directed rooftop with its two wedges.
\begin{tikzpicture}[line cap=round, line join=round, scale=0.95]
    \def\h{1.2}
    \foreach \i in {-2,...,1} { \draw[dashed] ({(\i+0.5)*\h},{-1.5*\h-0.4}) -- ({(\i+0.5)*\h},{1.5*\h+0.4}); }
    \foreach \j in {-2,...,1} { \draw[dashed] ({-1.5*\h-0.4},{(\j+0.5)*\h}) -- ({1.5*\h+0.4},{(\j+0.5)*\h}); }
    \foreach \i/\j in {-1/0,1/0,0/-1,0/1} { \fill[PairedA] ({(\i-0.5)*\h},{(\j-0.5)*\h}) rectangle ++(\h,\h); }
    \foreach \i/\j in {-1/-1,1/-1,-1/1,1/1} { \fill[PairedC] ({(\i-0.5)*\h},{(\j-0.5)*\h}) rectangle ++(\h,\h); }
    \fill[PairedE] ({-0.5*\h},{-0.5*\h}) rectangle ++(\h,\h);
    \foreach \i in {-1,...,2} { \draw[thick] ({(\i-0.5)*\h},{-1.5*\h}) -- ({(\i-0.5)*\h},{1.5*\h}); }
    \foreach \j in {-1,...,2} { \draw[thick] ({-1.5*\h},{(\j-0.5)*\h}) -- ({1.5*\h},{(\j-0.5)*\h}); }
    \draw[PairedF, line width=1.6pt] ({-0.5*\h},{-0.5*\h}) rectangle ++(\h,\h);
    \foreach \x/\y in {-0.5/-0.5,0.5/-0.5,-0.5/0.5,0.5/0.5} { \fill[PairedF] ({\x*\h},{\y*\h}) circle (2pt); }
    \foreach \i/\j/\l in {-1/-1/{-1,-1},0/-1/{0,-1},1/-1/{1,-1},-1/0/{-1,0},1/0/{1,0},-1/1/{-1,1},0/1/{0,1},1/1/{1,1}}
        { \node[font=\footnotesize] at ({\i*\h},{\j*\h}) {$(\l)$}; }
    \node at (0,0) {$\Om{\V{0}}$};
    \draw[thick,->] ({-1.5*\h-1.1},{-1.5*\h-0.4}) -- ++(0.6,0) node[anchor=north]{$\UV{x}$};
    \draw[thick,->] ({-1.5*\h-1.1},{-1.5*\h-0.4}) -- ++(0,0.6) node[anchor=east]{$\UV{y}$};
    \begin{scope}[shift={({-1.5*\h},{-1.5*\h-0.95})}, font=\small]
        \fill[PairedE] (0,0) rectangle ++(0.3,0.3); \node[anchor=west] at (0.35,0.15) {coincident};
        \fill[PairedA] (2.15,0) rectangle ++(0.3,0.3); \node[anchor=west] at (2.5,0.15) {edge};
        \fill[PairedC] (3.5,0) rectangle ++(0.3,0.3); \node[anchor=west] at (3.85,0.15) {vertex};
    \end{scope}
\end{tikzpicture}
\hspace{0.6cm}
\tdplotsetmaincoords{62}{20}
\begin{tikzpicture}[tdplot_main_coords, line cap=round, line join=round, scale=1.1]
    \def\dx{1.6}
    \def\dy{1.6}
    \def\roofh{1.5}
    \foreach \i in {0,...,2} { \draw[dashed] ({\i*\dx},0,0) -- ({\i*\dx},\dy,0); }
    \foreach \j in {0,...,1} { \draw[dashed] (0,{\j*\dy},0) -- (2*\dx,{\j*\dy},0); }
    \noteroofx{0}{0}{PairedF}{PairedE}{PairedB}{PairedA}
    \node at ({0.5*\dx},{0.5*\dy},0) {$\Om{mn}$};
    \node at ({1.5*\dx},{0.5*\dy},0) {$\Om{m+1,n}$};
    \node[PairedF] at ($(A)!0.5!(C)$) {$\xi$};
    \node[PairedB] at ($(B)!0.5!(F)$) {$1-\xi$};
    \node[anchor=west, font=\small] at ($(G)!0.6!(B)$) {$1$};
    \draw[thick,->] (-0.4,-1.2,0) -- (0.4,-1.2,0) node[anchor=north]{$\UV{x}$};
    \draw[thick,->] (-0.4,-1.2,0) -- (-0.4,-0.4,0) node[anchor=west]{$\UV{y}$};
    \draw[thick,->] (-0.4,-1.2,0) -- (-0.4,-1.2,0.8) node[anchor=south]{$\UV{z}$};
\end{tikzpicture}

    \caption{Left: the reference cell $\Om{\V{0}}$ and its eight touching neighbours $\Om{\offs}$, labelled by the offset
    $\offs\in\mathcal{T}$. The points with $R=0$ that the two cells have in common are drawn in red: all of
    $\Om{\V{0}}$ for the coincident cell, an edge for an edge neighbour and a vertex for a vertex neighbour. Right: the
    $x$-directed rooftop $\V{f}^x_{mn}$, split into its rising wedge $\xi$ on $\Om{mn}$ and its falling wedge $1-\xi$ on
    $\Om{m+1,n}$.}
    \label{fig:touching}
\end{figure}

\entryTitle{Galerkin entries}{galerkin} With~\eqref{eq:green}, the Galerkin entry of the \acs{EFIE} between rooftops
$\V{f}_p$ and $\V{f}_q$ is~\cite{machacek2025method}
\begin{equation}
    \label{eq:zpq}
    \begin{split}
    Z_{pq} = \frac{\I Z_0}{4\uppi k}\Bigl( & k^2 \int\!\!\int \V{f}_p(\V{r})\cdot\V{f}_q(\V{r}')\,g(R)\D A'\D A\\
    & - \int\!\!\int \nabla\cdot\V{f}_p(\V{r})\,\nabla'\cdot\V{f}_q(\V{r}')\,g(R)\D A'\D A \Bigr).
    \end{split}
\end{equation}
By translation invariance the entries depend only on the offset of the two rooftops. They are generated by
\begin{equation*}
    A^{xx}(\offs) = \int\!\!\int \V{f}^x_{\V{0}}\cdot\V{f}^x_{\offs}\,g, \qquad
    A^{yy}(\offs) = \int\!\!\int \V{f}^y_{\V{0}}\cdot\V{f}^y_{\offs}\,g, \qquad
    \Phi(\offs)  = \int_{\Om{\V{0}}}\!\int_{\Om{\offs}} g .
\end{equation*}
The divergences are piecewise constant, so the scalar-potential parts of all blocks are finite differences of $\Phi$; for
example
$\int\!\!\int \nabla\cdot\V{f}^x_{\V{0}}\,\nabla'\cdot\V{f}^x_{\offs}\,g
= a^{-2}\bigl(2\Phi(\offs)-\Phi(\offs+\ex)-\Phi(\offs-\ex)\bigr)$.
This is how \src{} assembles the three Toeplitz generators.

\defTitle{Cell-pair moments}{moments} Let $K:(0,\infty)\to\C$ satisfy $K(\NORM{\cdot})\in L^1_{\mathrm{loc}}(\Rn^2)$. For
$\offs\in\Zn^2$ and $\alpha,\beta\in\{0,1\}$ define
\begin{align*}
    \MomK{K}{\alpha\beta}{\offs} &= \frac{1}{(ab)^2}\int_{\Om{\V{0}}}\int_{\Om{\offs}}
        \xi(\V{r})^{\alpha}\,\xi(\V{r}')^{\beta}\,K\bigl(\NORM{\V{r}-\V{r}'}\bigr)\D A(\V{r}')\D A(\V{r}), \\
    N^{[K]}_{\alpha\beta}(\offs) &= \frac{1}{(ab)^2}\int_{\Om{\V{0}}}\int_{\Om{\offs}}
        \eta(\V{r})^{\alpha}\,\eta(\V{r}')^{\beta}\,K\bigl(\NORM{\V{r}-\V{r}'}\bigr)\D A(\V{r}')\D A(\V{r}),
\end{align*}
where $\xi(\V{r}')$ and $\eta(\V{r}')$ are the local coordinates in $\Om{\offs}$. For $K(R)=R^p$ we write $\Mom{p}{\alpha\beta}{\offs}$
and $\Nom{p}{\alpha\beta}{\offs}$. Since $N_{00}=M_{00}$, there are seven moments, ordered in \src{} as $M_{00}$, $M_{10}$, $M_{01}$,
$M_{11}$, $N_{10}$, $N_{01}$, $N_{11}$.

\vspace{\baselineskip}
The moments carry all the geometry. The next lemma lists their symmetries.

\lemmaTitle{Symmetries of the moments}{symmetries} Let $K$ be admissible, $\offs=(d_x,d_y)\in\Zn^2$ and
$\alpha,\beta\in\{0,1\}$. Then:
\begin{enumerate}[(i)]
    \item $\MomK{K}{\alpha\beta}{d_x,-d_y} = \MomK{K}{\alpha\beta}{d_x,d_y}$;
    \item reflection in $x$:
    \begin{align*}
        \MomK{K}{00}{-d_x,d_y} &= \MomK{K}{00}{\offs}, \\
        \MomK{K}{10}{-d_x,d_y} &= \bigl(M^{[K]}_{00}-M^{[K]}_{10}\bigr)(\offs),\\
        \MomK{K}{01}{-d_x,d_y} &= \bigl(M^{[K]}_{00}-M^{[K]}_{01}\bigr)(\offs), \\
        \MomK{K}{11}{-d_x,d_y} &= \bigl(M^{[K]}_{00}-M^{[K]}_{10}-M^{[K]}_{01}+M^{[K]}_{11}\bigr)(\offs);
    \end{align*}
    \item $N^{[K]}_{\alpha\beta}(d_x,d_y;a,b) = \MomK{K}{\alpha\beta}{d_y,d_x;b,a}$ (the cell sizes written explicitly);
    \item $\MomK{K}{\alpha\beta}{\offs} = \MomK{K}{\beta\alpha}{-\offs}$.
\end{enumerate}

\vspace{\baselineskip}
\proof
The kernel depends only on $\NORM{\V{r}-\V{r}'}$, which is invariant under reflections, translations and the interchange
of $\V{r}$ and $\V{r}'$.

(i) The reflection $y\mapsto b-y$ maps $\Om{\V{0}}$ onto itself and $\Om{(d_x,-d_y)}$ onto $\Om{(d_x,d_y)}$. It leaves
both local $x$-coordinates unchanged.

(ii) The reflection $x\mapsto a-x$ maps $\Om{\V{0}}$ onto itself with $\xi\mapsto1-\xi$, and $\Om{(-d_x,d_y)}$ onto
$\Om{(d_x,d_y)}$ with $\xi'\mapsto1-\xi'$. Indeed, $[-d_xa,(1-d_x)a]$ goes to $[d_xa,(d_x+1)a]$ and
$x'/a+d_x\mapsto 1-(x'/a+d_x)$. Expand $(1-\xi)^\alpha(1-\xi')^\beta$.

(iii) Exchange the roles of $x$ and $y$.

(iv) Interchange $\V{r}\leftrightarrow\V{r}'$ and translate by $-\offs$; the cell $\Om{\offs}$ becomes $\Om{\V{0}}$ and
$\Om{\V{0}}$ becomes $\Om{-\offs}$.

\qed

\propTitle{Generators from moments}{generators} For $\offs=(d_x,d_y)$ with $d_x\ge0$,
\begin{equation}
    \label{eq:axx}
    \begin{split}
    \frac{A^{xx}(\offs)}{(ab)^2} = {}& M_{11}(\offs) + \bigl(M_{00}-M_{10}-M_{01}+M_{11}\bigr)(\offs)\\
    &+ \bigl(M_{10}-M_{11}\bigr)(\offs+\ex) + \bigl(M_{01}-M_{11}\bigr)(\offs-\ex),
    \end{split}
\end{equation}
with $M=M^{[g]}$. For $d_x=0$ the last term is to be read as
\begin{equation*}
    \bigl(M_{01}-M_{11}\bigr)(-1,d_y) = \bigl(M_{10}-M_{11}\bigr)(1,d_y).
\end{equation*}
Furthermore
$\Phi(\offs)=(ab)^2M^{[g]}_{00}(\offs)$, and $A^{yy}$ is given by~\eqref{eq:axx} with $N$ and $\ey$ in place of $M$ and $\ex$.

\vspace{\baselineskip}
\proof
Split both rooftops into their wedges, as in Definition~\defRef{rooftop}. The reference rooftop $\V{f}^x_{\V{0}}$ has the
rising wedge $\xi$ on $\Om{\V{0}}$ and the falling wedge $1-\xi$ on $\Om{\ex}$. The rooftop $\V{f}^x_{\offs}$ has the
rising wedge $\xi'$ on $\Om{\offs}$ and the falling wedge $1-\xi'$ on $\Om{\offs+\ex}$. Both are $x$-directed, so the dot
product is the product of the profiles. The four wedge pairs of Figure~\ref{fig:wedgepairs} give
\begin{center}
    \renewcommand{\arraystretch}{1.25}
    \setlength{\tabcolsep}{5pt}
    \begin{tabular}{llccc}
        \toprule
        pair & cells & offset & weight & moment \\ \midrule
        1, rising--rising     & $\Om{\V{0}}$, $\Om{\offs}$     & $\offs$     & $\xi\xi'$           & $M_{11}$ \\
        2, falling--falling   & $\Om{\ex}$, $\Om{\offs+\ex}$   & $\offs$     & $(1-\xi)(1-\xi')$   & $M_{00}-M_{10}-M_{01}+M_{11}$ \\
        3, rising--falling    & $\Om{\V{0}}$, $\Om{\offs+\ex}$ & $\offs+\ex$ & $\xi(1-\xi')$       & $M_{10}-M_{11}$ \\
        4, falling--rising    & $\Om{\ex}$, $\Om{\offs}$       & $\offs-\ex$ & $(1-\xi)\xi'$       & $M_{01}-M_{11}$ \\ \bottomrule
    \end{tabular}
\end{center}
with the moments taken at the offset of the pair (for pair 4 after the translation by $-\ex$).
The normalization $(ab)^{-2}$ in Definition~\defRef{moments} accounts for the factor $(ab)^2$. For $d_x=0$ the offset
$\offs-\ex=(-1,d_y)$ is negative. Lemma~\lemmaRef{symmetries}(ii) gives
\begin{align*}
    (M_{01}-M_{11})(-1,d_y) &= (M_{00}-M_{01})(1,d_y)-(M_{00}-M_{10}-M_{01}+M_{11})(1,d_y)\\
                            &= (M_{10}-M_{11})(1,d_y).
\end{align*}
The statement on $\Phi$ is the definition with $\alpha=\beta=0$. The $y$-direction is identical.

\qed

\begin{figure}[ht]
    \centering
    % The four wedge pairs of two x-directed rooftops at the offset d = (2,2) (Proposition generators).
\tdplotsetmaincoords{62}{20}
\begin{tikzpicture}[tdplot_main_coords, line cap=round, line join=round, scale=1.35]
    \def\dx{1.3}
    \def\dy{1.3}
    \def\roofh{1.1}
    \foreach \i in {0,...,4} { \draw[dashed] ({\i*\dx},0,0) -- ({\i*\dx},3*\dy,0); }
    \foreach \j in {0,...,3} { \draw[dashed] (0,{\j*\dy},0) -- (4*\dx,{\j*\dy},0); }
    \draw[thick, ->, >=stealth] (0,0,0) -- ({2*\dx},{2*\dy},0) node[pos=0.8, anchor=north west]{$\offs$};
    \noteroofx{0}{0}{PairedF}{PairedE}{PairedF}{PairedE}
    \node[font=\small] at ({0.5*\dx},{0.5*\dy},0) {$\Om{\V{0}}$};
    \node[font=\small] at ({1.5*\dx},{0.5*\dy},0) {$\Om{\ex}$};
    \noteroofx{2}{2}{PairedB}{PairedA}{PairedB}{PairedA}
    \node[font=\small] at ({2.5*\dx},{2.5*\dy},0) {$\Om{\offs}$};
    \node[font=\small] at ({3.5*\dx},{2.5*\dy},0) {$\Om{\offs+\ex}$};
    \coordinate (r0) at ({0.6*\dx},{0.5*\dy},{0.6*\roofh});
    \coordinate (f0) at ({1.4*\dx},{0.5*\dy},{0.6*\roofh});
    \coordinate (rd) at ({2.6*\dx},{2.5*\dy},{0.6*\roofh});
    \coordinate (fd) at ({3.4*\dx},{2.5*\dy},{0.6*\roofh});
    \tikzset{pair/.style={dotted, very thick, <->, >=stealth}, num/.style={solid, circle, draw, fill=white, inner sep=1pt, font=\footnotesize}}
    \draw[pair] (r0) .. controls +(0,0,1.6) and +(0,0,1.6) .. node[num, pos=0.35]{1} (rd);
    \draw[pair] (f0) .. controls +(0,0,2.6) and +(0,0,2.6) .. node[num, pos=0.62]{2} (fd);
    \draw[pair] (r0) .. controls +(0,0,3.4) and +(0,0,3.4) .. node[num, pos=0.5]{3} (fd);
    \draw[pair] (f0) .. controls +(0,0,0.8) and +(0,0,0.8) .. node[num, pos=0.5]{4} (rd);
    \node[PairedF, anchor=east] at ({0,0,0.4*\roofh}) {$\V{f}^x_{\V{0}}$};
    \node[PairedB, anchor=west] at ({4*\dx},{3*\dy},{0.3*\roofh}) {$\V{f}^x_{\offs}$};
\end{tikzpicture}

    \caption{The four wedge pairs of the rooftops $\V{f}^x_{\V{0}}$ and $\V{f}^x_{\offs}$ in
    Proposition~\propRef{generators}, drawn for $\offs=(2,2)$, numbered as in its proof. Every pair is a pair of cells
    at one offset, weighted by a product of two linear profiles.}
    \label{fig:wedgepairs}
\end{figure}

\vspace{\baselineskip}
The generators of all three Toeplitz blocks need the seven moments on the offsets
$d_x\ge-1$, $d_y\ge0$, and by Lemma~\lemmaRef{symmetries} on $d_x,d_y\ge0$ only.
